\section{La ruta de evaluación de AI for Math: del banco de problemas a la contribución a la investigación}

Para saber si AI for Math avanza de verdad, no basta con contar cuántos problemas de competencia resuelve el modelo. Los problemas de competencia tienen valor porque sus respuestas son claras y su retroalimentación es estricta, pero son solo el punto de entrada. Una evaluación de nivel superior debería considerar la tasa de aprobación de las demostraciones formales, las contribuciones a la biblioteca de teoremas, el descubrimiento de lemmas, la construcción de contraejemplos, la eficiencia de la colaboración con matemáticos y la aparición de nuevas líneas de investigación.

\subsection{Niveles de evaluación}

\noindent\begin{longtable}{p{0.22\textwidth}p{0.35\textwidth}p{0.34\textwidth}}
\toprule
Nivel & Indicadores medibles & Limitaciones \\
\midrule
Resolución de problemas de competencia & Tasa de aciertos, pass@k, tiempo de resolución & Es fácil sobreajustarse al banco de problemas; no refleja capacidad de investigación. \\
Completado formal & Tasa de aprobación de Lean theorem, tasa de éxito de tactic & Depende de la cobertura de la biblioteca de teoremas y de la selección de problemas. \\
Generación de demostraciones & Longitud de la demostración completa, legibilidad, tasa de aprobación de la verificación & Correcta no equivale a elegante; una demostración larga puede no aportar ninguna idea de fondo. \\
Contribución a la biblioteca de teoremas & Número de lemmas nuevos, tasa de reutilización, costo de mantenimiento & Requiere la aceptación de la comunidad; no basta con que la máquina los genere. \\
Colaboración en investigación & Tiempo que ahorran los matemáticos, contraejemplos encontrados, conjeturas propuestas & Es difícil de estandarizar, pero es lo más cercano al valor real. \\
\bottomrule
\end{longtable}

\begin{importantbox}{El núcleo de la ruta de evaluación}
La evaluación de AI for Math debería pasar gradualmente de «la respuesta es correcta» a «la demostración es verificable», «la demostración tiene poder explicativo», «la biblioteca de teoremas es reutilizable» y «los matemáticos están dispuestos a colaborar». Cuanto más se avanza, más difícil es cuantificar, pero también más cerca se está del valor real.
\end{importantbox}

\subsection{Por qué esto se relaciona con los Agents}

AI for Math parece un problema científico, pero en realidad también es un problema de Agents: el modelo tiene que entrar en el entorno de Lean, leer el goal state, invocar tactics, observar los errores, corregir la demostración y seguir avanzando. Este ciclo cerrado se parece mucho al de un web agent o un coding agent; la diferencia es que el entorno es más formal y la retroalimentación más estricta. Por eso, EP137 puede verse como una instancia científica de la historia técnica de los Agents que aborda EP139.

\begin{knowledgebox}{El entorno de Lean es un entorno de Agents de alta calidad}
Lean ofrece estados definidos, acciones definidas, verificación definida y retroalimentación de errores definida. Comparado con un entorno web abierto, es más riguroso; comparado con una conversación ordinaria, es más fácil de entrenar. Esto convierte a AI for Math en un campo de pruebas ideal para estudiar la capacidad de aprendizaje de los Agents.
\end{knowledgebox}

\subsection{Resumen de la sección}

La evaluación de AI for Math no puede quedarse en la tasa de aciertos sobre un banco de problemas. La ruta real avanza desde la resolución de problemas, la formalización y la generación de demostraciones hacia la contribución a la biblioteca de teoremas y la colaboración con matemáticos.

